The three-epoch full-validation experiment followed a smaller adaptation screen, which should
remain visible as completed work. It used 64 selected training frames
(63 effective after filtering), seed 20260823, three RPN epochs and
three RCNN epochs. Six input conditions were evaluated before and after
adaptation on the same 256 validation frames; all twelve evaluations
completed. Table~\ref{tab:k-screen} reports their Moderate results under
the explicitly recorded legacy Python evaluator. Change is adapted
minus frozen AP; Gap is adapted AP minus the matching adapted baseline,
both in AP points. These values are not points on an R40 learning curve.
The complete Easy/Moderate/Hard triplets are retained in the archived
experiment outputs.

\begin{table}[!t]\centering\setlength{\tabcolsep}{4pt}
\begin{tabular}{llrrrr}\toprule
Line & Input & Frozen & Adapted & Change & Gap\\\midrule
A & Original & 79.19 & 78.18 & $-1.01$ & 0.00\\
A & PU-GCN & 60.27 & 64.44 & $+4.17$ & $-13.74$\\
A & PDANS & 68.58 & 67.02 & $-1.57$ & $-11.17$\\
B & Sparse & 66.62 & 62.91 & $-3.70$ & 0.00\\
B & PU-GCN & 32.92 & 37.52 & $+4.59$ & $-25.40$\\
B & PDANS & 43.01 & 47.43 & $+4.43$ & $-15.48$\\\bottomrule
\end{tabular}
\caption{PointRCNN adaptation screen on 256 frames: Car Moderate 3D AP
at IoU 0.70 under the legacy evaluator.}\label{tab:k-screen}
\end{table}

The screen recovered more than four AP points for both PU-GCN inputs
and for sparse PDANS, but not for dense PDANS. Both adapted baselines
also deteriorated. Thus a positive before/after change alone would
overstate the benefit: all adapted generated inputs remained below
their paired adapted baseline. The later stage expanded training to
3,712 frames, initially evaluated a 256-frame pilot, and ultimately
completed the 3,769-frame, 20-arm matrix. Increasing training coverage,
changing the integration configuration and changing the evaluation
scope prevent these stages from being interpreted as a single
controlled sample-size ablation.
